\documentclass[11pt]{article}
% Shared preamble for all daily exercises
% Update this file to change settings (padding, parindent, etc.) for all exercises

\usepackage[margin=0.75in]{geometry}
\usepackage{amsmath}
\usepackage{amssymb}

% Custom settings
\setlength{\parindent}{0pt}
\setlength{\parskip}{0.2cm}

\title{Daily Exercise}
\author{Dhruman Gupta}
\date{5th March}

\begin{document}

% Common header for daily exercises
% This file can be included in other LaTeX documents

\maketitle

\noindent
\textbf{Course:} CS-3330, Theory of Computation

\vspace{0.2cm}

\noindent
\textbf{Question:}\noindent 
Construct a Turing machine that takes as input a number n in unary and accepts if and only if n is a perfect square.

\textbf{Answer:}

Let $M$ be a turing machine. The initial configuration is $(q_0, 1^n\sqcup, 0)$. The machine will then do the following:

\begin{itemize}
    \item Read until the first $\sqcup$.
    \item Add a $\#$ as a separator.
    \item Initialise $i = 1$.
    \item While there exists a $1$ on the left of $\#$:
    \begin{itemize}
        \item Write $i$ many $1$ to the right of $\#$.
        \item While there exists a $1$ on the right of $\#$:
        \begin{itemize}
            \item Swap the right-most $1$ with empty, and replace the right-most $1$ on the left of $\#$ with a $X$.
            \item If no $1$ exists on the left of $\#$, then reject.
        \end{itemize}
        \item If all symbols on the left of $\#$ are $X$, then accept.
        \item Increment $i$ by 2.
    \end{itemize}
\end{itemize}

This is a valid Turing machine that accepts if and only if n is a perfect square. This stems from the identity that every perfect square is a sum of consecutive odd numbers.
\end{document}
